\documentclass{article}
\usepackage[T1]{fontenc}
\usepackage{lmodern}
\usepackage{graphicx}
\usepackage{amsmath,amssymb}
\usepackage[a4paper,margin=2cm]{geometry}
\usepackage[absolute,overlay]{textpos}
\setlength{\TPHorizModule}{1mm}
\setlength{\TPVertModule}{1mm}
\usepackage[indonesian]{babel}
\usepackage{hyperref}
\setlength{\emergencystretch}{2em}
% unit: Halaman 14

\begin{document}

% === Halaman 14 (llm) ===

\begin{itemize}
    \item[$\bullet$] \textbf{OTP ini tidak dapat dipecahkan karena:}
    \begin{enumerate}
        \item Kunci OTP (acak) + plainteks (tidak acak) = cipherteks (acak)
        
        \begin{align*}
            \text{Enkripsi: } C &= (P + K) \pmod{26} \\ 
            \text{Dekripsi: } P &= (C - K) \pmod{26}
        \end{align*}
        
        \item Hanya terdapat satu kunci yang memetakan plainteks ke cipherteks, begitu juga sebaliknya
        \begin{itemize}
            \item[$\rightarrow$] fungsi $C = (P + K) \pmod{26}$ adalah pemetaan satu-ke-satu (\textit{one-to-one}).
            \item[$\rightarrow$] artinya, untuk sembarang plainteks dan cipherteks hanya ada satu kunci yang memetakannya satu sama lain
            \item[$\rightarrow$] Akibatnya, jika kriptanalis mencoba mendekripsi cipherteks dengan beberapa kunci berbeda ada kemungkinan menghasilkan beberapa plainteks yang bermakna, sehingga kriptanalis kesulitan menentukan plainteks mana yang benar.
        \end{itemize}
    \end{enumerate}
\end{itemize}

\end{document}
